\chapter{Marco Metodológico y Diseño de Investigación}
\label{ch:metodologia}

Este capítulo detalla y justifica las elecciones metodológicas, el tipo de diseño, la matriz operacional de variables y las técnicas econométricas que conforman el protocolo de esta investigación.

\section{Enfoque, Tipo y Alcance de la Investigación}
Siguiendo la clasificación estándar de los diseños de investigación en ciencias sociales \citep{marradi2007, rojas2013} y los criterios de formulación de un proyecto de tesis \citep{bassi2015}, esta investigación adopta un \textbf{enfoque cuantitativo} y un \textbf{alcance descriptivo-\allowbreak{}explicativo}, buscando establecer relaciones estadísticas e inferir umbrales críticos de reacción fiscal. El diseño temporal es \textbf{longitudinal de series de tiempo} (caso único, Argentina), lo cual permite capturar la volatilidad e inestabilidad macroeconómica local sin el sesgo de promedios multinacionales. Siguiendo una \textbf{lógica deductiva}, se derivan hipótesis contrastables ($H_{1a}$, $H_{1b}$, $H_{1c}$) desde la Restricción Presupuestaria Intertemporal y se someten a contrastación empírica en el Capítulo \ref{ch:resultados}.

La investigación trabaja con \textbf{dos ventanas muestrales complementarias}, no una única. La ventana original, 2004T1--2025T4 ($n=88$), cubre el período con cobertura primaria homogénea de todas las variables del modelo y sostiene la estimación DOLS de referencia (Sección \ref{sec:dols}) junto con el resto del protocolo de robustez (umbral de Hansen, quiebres de Bai-Perron, DSA estocástico). La ventana ampliada, 1999T1--2025T4 ($n=108$), incorpora un tramo empalmado 1996--2003 (Sección \ref{sec:empalme}) y sostiene la estimación VECM/VAR restringido que reemplaza a DOLS como técnica de referencia para la Función de Reacción Fiscal de largo plazo (Sección \ref{sec:vecm}), a pedido del director de tesis, quien solicitó explícitamente extender la ventana temporal aun a costa de incorporar quiebres estructurales extremos (el colapso de la convertibilidad y la cesación de pagos de 2001--2002), exigiendo en contrapartida que el procedimiento de extensión quede documentado en su totalidad.

\subsection{Unidad de Análisis}
\label{sec:unidad_analisis}
La Unidad de Análisis es el \textbf{trimestre macroeconómico consolidado} de Argentina, para el período 2004T1--2025T4 en la ventana original y 1999T1--2025T4 en la ventana ampliada. Esta frecuencia, frente a la agregación anual tradicional, captura la volatilidad de los flujos fiscales y las oscilaciones del riesgo país sin los suavizamientos de la baja frecuencia anual. El trimestre es también el límite superior de desagregación posible: el Resultado Primario y la Deuda Bruta del SPNF consolidado --a diferencia de los datos de Administración Nacional-- no tienen publicación mensual verificable en fuente primaria para toda la ventana muestral.

\subsection{Matriz Operacional de Variables}
\label{sec:matriz_operacional}
Siguiendo la estructura del dato científico --unidad de análisis, variable, valor e indicador-- propuesta por \citet{samaja1993}, la Tabla \ref{tab:variables_operacional} formaliza la correspondencia entre cada dimensión teórica desarrollada en el Capítulo \ref{ch:marco_teorico}, su variable econométrica simbólica y su indicador operacional efectivamente medido, explicitando además su unidad de medida y su fuente de datos oficial.

\begin{table}[H]
\centering
\caption{\label{tab:variables_operacional} Matriz de Operacionalización de Variables e Isomorfismo Lógico.}
\resizebox{\textwidth}{!}{%
\footnotesize
\begin{tabular}{p{2.2cm}p{1.1cm}p{3.6cm}p{1.5cm}p{3.4cm}p{2.2cm}}
\toprule
\textbf{Dimensión Teórica} & \textbf{Variable} & \textbf{Indicador Operacional / de Aproximación} & \textbf{Unidad} & \textbf{Transformación Aplicada} & \textbf{Fuente} \\
\midrule
Esfuerzo Fiscal Primario Corriente & $pb_t$ & Resultado Primario del SPNF / PIB & \% & Nivel, sin transformar (admite negativos) & Hacienda (MECON) \\
Stock de Deuda Heredado & $d_{t-1}$ & Stock de Deuda Bruta Consolidada SPNF$_{t-1}$ / PIB & \% & Nivel, rezagado un período & Finanzas (MECON) \\
Ciclo Macroeconómico & $\tilde{y}_t$ & Brecha del Producto (HP s/ PIB real) & Desv. \% & Filtro HP ($\lambda=1600$), \% de desvío tendencial & INDEC \\
Perturbaciones de Gasto Extraordinario & $g_{t}^{exp}$ & Gasto primario en crisis real & p.p. PIB & Nivel, sin transformar & MECON \\
Restricción Financiera Internacional & $risk_t$ & Riesgo País (EMBI+) & pb & Nivel; rezagado como instrumento (Etapas 3 y 5) o contemporáneo como variable endógena (Etapa 4) & BCRA / Bloomberg \\
Efecto de Valoración Cambiaria & $tc_t^{val}$ & Interacción deuda en moneda extranjera $\times$ variación TCRM & Tasa var. & Tasa de variación interperiódica del TCRM & BCRA / INDEC \\
Tensión Cambiaria (nivel) & --- & TCR \textbf{Multilateral} (TCRM, base BCRA; no bilateral USD/ARS) & Índice & Nivel; variable de control macroeconómico & BCRA / datos.gob.ar \\
Precios Relativos & --- & Coef. de Estabilización de Referencia & Índice acum. & Nivel, sin transformar & BCRA \\
Perturbación Financiera / Riesgo Emergente & --- & VIX (CBOE) y ETF EMB (iShares JPMorgan EMBI) & Puntos / USD & Nivel; instrumentos externos (IV-2SLS) & Yahoo Finance / Bloomberg \\
\midrule
Restricción Externa (paradigma marco, Nivel I) & --- & \textit{Sin variable de aproximación propia.} Paradigma (Cap. \ref{ch:marco_teorico}) operacionalizado vía $risk_t$ y $tc_t^{val}$ & N/A & N/A & --- \\
\bottomrule
\end{tabular}%
}
\begin{flushleft}
\footnotesize \textit{Nota.} Elaboración propia en base al libro de códigos de la matriz de datos consolidada del proyecto. La Deuda Pública Consolidada corresponde exclusivamente al Sector Público No Financiero (SPNF), excluyendo pasivos remunerados del BCRA (ver aclaración de perímetro institucional a continuación). Fuentes institucionales: Ministerio de Economía --- \url{https://www.economia.gob.ar} ---; INDEC --- \url{https://www.indec.gob.ar} ---; BCRA --- \url{https://www.bcra.gob.ar} ---; Portal de Datos Abiertos --- \url{https://www.datos.gob.ar} ---; Yahoo Finance --- \url{https://finance.yahoo.com} ---.
\end{flushleft}
\end{table}

La variable Deuda Pública ($d_t$) corresponde al stock bruto consolidado del Sector Público No Financiero (SPNF), según los estándares del Manual de Estadísticas de Finanzas Públicas del FMI (2014) y el perímetro institucional que fija la Ley de Administración Financiera \citep{ley24156} para la Secretaría de Finanzas (MECON). Se excluyen de la definición central los pasivos remunerados del BCRA (LEBAC, LELIQ y pases), instrumentos de política monetaria emitidos al amparo de las facultades que le confiere su Carta Orgánica \citep{ley24144}, cuya dinámica responde a objetivos distintos de los del financiamiento fiscal directo. No obstante, como prueba de robustez (Capítulo \ref{ch:resultados}), se construye una serie consolidada alternativa que suma los pasivos del BCRA al stock del SPNF.

\section{Construcción y Crítica de las Fuentes de Datos}
El trabajo empírico en Argentina enfrenta el desafío estructural de la calidad y continuidad de las estadísticas públicas. En línea con la exigencia de respaldo empírico de toda afirmación sobre la evolución de las variables \citep{ynoub2007}, esta investigación utiliza exclusivamente \textbf{fuentes secundarias oficiales}, obtenidas mediante un protocolo de ingesta en dos etapas. En primera instancia, se intentó la descarga automatizada vía \textit{Application Programming Interface} (API): la Deuda Pública, el Resultado Primario y el Tipo de Cambio Real Multilateral desde \texttt{apis.datos.gob.ar} (MECON/BCRA/INDEC); el Coeficiente de Estabilización de Referencia desde la API pública del BCRA; y el índice de volatilidad VIX desde Yahoo Finance. En segunda instancia, ante discontinuidades o interrupciones de conectividad, se recurrió a un esquema jerárquico de fuentes de contingencia (informes del BCRA y MECON, series históricas de JP Morgan y CEPAL, y la base \textit{World Economic Outlook} del FMI), exigiendo en todos los casos el cruce contra un mínimo de dos fuentes independientes antes de su incorporación a la base consolidada, como criterio de confiabilidad instrumental de la matriz de datos \citep{cohen2019}. Para el Resultado Primario y la Deuda Pública en particular, los identificadores de serie de \texttt{apis.datos.gob.ar} vigentes al inicio del proyecto quedaron discontinuados durante su desarrollo; en la práctica, la totalidad de las 88 observaciones trimestrales de ambas variables proviene de esta segunda vía, sin un llamado API exitoso de por medio, consistente entre sí y con la base FMI WEO en los puntos de control verificados.

\subsection{Empalme de la Ventana Ampliada (1996--2003)}
\label{sec:empalme}
Ninguna de las cuatro variables del sistema de cointegración (deuda pública, resultado primario, EMBI+, TCRM) tiene cobertura trimestral verificable en fuente primaria antes de 2004 o, en el caso del EMBI+, antes de 1998. Extender la ventana muestral a 1999T1--2025T4 exigió, entonces, construir el tramo 1996--2003 mediante un procedimiento de empalme documentado en su totalidad, y no una simple extrapolación.

Para la Deuda Pública y el Resultado Primario del SPNF, la fuente del tramo faltante es el \textit{Compendio Fiscal 1993--2006} de la Secretaría de Hacienda (MECON), de frecuencia anual. El procedimiento consta de dos pasos, aplicados por separado a cada variable. Primero, dado que las series del Compendio no miden exactamente lo mismo que las series API vigentes desde 2004 (distinto perímetro institucional y, en el caso de la deuda, distinta valuación), se calcula un \textbf{coeficiente de enlace} en el año de solapamiento más cercano al punto de empalme (2004): la razón entre el valor de la serie moderna y el valor de la serie del Compendio en ese año. Ese coeficiente, que no supone identidad entre ambas series, es el que se aplica a los valores anuales históricos antes de desagregarlos. El coeficiente resultó razonablemente estable para la deuda ($\approx 0.87$) pero inestable para el resultado primario, cuyo cociente Compendio/dataset-actual pasa de 1.12 en 2004 a 2.07 en 2006, lo cual sugiere una diferencia de perímetro más profunda que una simple diferencia de valuación; se utiliza igualmente el coeficiente del año de empalme por ser el más cercano al punto de unión, y esta limitación queda declarada como la aproximación menos confiable del procedimiento. Segundo, los valores anuales reescalados se desagregan a frecuencia trimestral mediante el método de Denton sin indicador \citep{denton1971}, que minimiza la suma de las diferencias trimestre a trimestre al cuadrado sujeto a que el promedio de los cuatro trimestres de cada año reproduzca el valor anual dado. Este es el procedimiento estándar para pasar de frecuencia anual a trimestral cuando no existe una serie de alta frecuencia relacionada que sirva de indicador. Para evitar un salto artificial en el punto de unión, la optimización incluye como puntos fijos (sin restricción de promedio anual, solo como ancla del término de suavidad) los trimestres reales 2004T1--2006T4 inmediatamente posteriores al tramo estimado.

Para el PIB real, al no existir la ambigüedad de perímetro institucional del resultado fiscal, el nivel anual 1996--2003 se reconstruye encadenando hacia atrás desde el propio dato trimestral real de 2004 con las tasas de crecimiento anuales publicadas por INDEC y el Banco Mundial, sin coeficiente de enlace; la desagregación trimestral sigue el mismo procedimiento de Denton anclado contra 2004--2006.

Para el TCRM y el EMBI+, en cambio, no se empalma: en la matriz de datos ampliada, se reemplaza la totalidad del panel (incluida la ventana 2004--2025) por series reales de fuente primaria, en lugar de heredar de la matriz original los valores de contingencia usados hasta este punto. El TCRM proviene del Índice de Tipo de Cambio Real Multilateral publicado por el BCRA (rebasado a diciembre de 2001 $=1$, cobertura 1997--2025), que reemplaza al valor de contingencia utilizado hasta este punto; se trata de un hallazgo de auditoría de fuentes, documentado en la Sección \ref{sec:auditoria_fallback}. El EMBI+ (riesgo país) proviene de la serie diaria publicada por Ámbito Financiero, agregada a frecuencia trimestral por promedio simple, con cobertura desde diciembre de 1998; este es el límite que fija el inicio efectivo de la ventana ampliada en 1999T1, un trimestre después del primer dato disponible. El VIX y el CER permanecen en su ventana original (2004--2025, 88 observaciones): no forman parte del sistema de cointegración, pues son variables de control o instrumentos externos, y su extensión no es necesaria para la estimación VECM.

La columna \texttt{es\_interpolado} de la matriz de datos ampliada distingue explícitamente, para cada trimestre, si el dato proviene del empalme (1996T1--2003T4, deuda/resultado primario/PIB real únicamente) o de fuente primaria directa, preservando la trazabilidad exigida por el protocolo de fuentes secundarias \citep{cohen2019}.

\subsection{Auditoría de Cobertura: Series de Contingencia vs. Fuente Primaria}
\label{sec:auditoria_fallback}
Un chequeo de trazabilidad, consistente en comparar cada columna de la matriz de datos original contra los valores de contingencia embebidos en el código de ingesta, reveló que, además de la Deuda Pública y el Resultado Primario (ya documentados como no provenientes de un llamado API exitoso), el TCRM y el EMBI+ también correspondían en su totalidad, en las 88 observaciones originales, a la fuente de contingencia y no a una consulta API exitosa. Esta auditoría motivó directamente su reemplazo por las series reales BCRA y Ámbito Financiero descritas en la Sección \ref{sec:empalme} en la matriz de datos ampliada, aplicado no solo en el tramo empalmado sino en la totalidad de su ventana 1999--2025, verificado mediante el cruce en el subperíodo de solapamiento 2004--2025 contra los valores previamente utilizados. La matriz de datos original de 88 observaciones, que sostiene DOLS, IV-2SLS, el umbral de Hansen, Bai-Perron y el DSA estocástico, conserva por el momento los valores de contingencia; la Sección \ref{sec:limitaciones_muestra} del Capítulo \ref{ch:descriptivo} y el Capítulo \ref{ch:resultados} documentan si y cómo se propaga este reemplazo a esos ejercicios.

\section{Estrategia Econométrica}

La evaluación de la fatiga fiscal requiere un abordaje econométrico secuencial de seis etapas, diseñado para evitar estimadores sesgados o regresiones espurias.

\subsection{Etapa 1: Estacionariedad y Quiebre Estructural Endógeno}
Se aplican los contrastes de Dickey-Fuller Aumentado (ADF) y Kwiatkowski-Phillips-Schmidt-Shin (KPSS) ---de hipótesis nulas opuestas--- junto con el test de \textit{Dickey-Fuller Generalized Least Squares} (DF-GLS) de \textcite{elliott1996} sobre cada serie en niveles y en primeras diferencias. Para los casos de diagnóstico ambiguo entre dichos contrastes, se implementa adicionalmente el test de quiebre estructural endógeno de \textcite{zivot1992}, que no requiere especificar \textit{a priori} la fecha de ruptura sino que la estima mediante búsqueda secuencial del punto que minimiza el estadístico $t$ del parámetro autorregresivo, permitiendo distinguir series genuinamente integradas de series estacionarias en torno a una tendencia quebrada. El protocolo se aplica sobre la ventana original (88 obs.) para la totalidad de las variables de la matriz de datos, y se repite sobre la ventana ampliada (108 obs.) para las cuatro variables del sistema de cointegración, con resultados de orden de integración consistentes entre ambas ventanas (Sección \ref{sec:estacionariedad}).

\subsection{Etapa 2: Selección de Rezagos y Cointegración de Johansen}
\label{sec:etapa2_johansen}
Antes de estimar cualquier relación de largo plazo, se verifica mediante el procedimiento de máxima verosimilitud de Johansen (contraste de Traza) si el sistema conformado por la ratio de deuda ($d_t$), el resultado primario ($pb_t$), el riesgo soberano ($risk_t$) y el tipo de cambio real ($tc_t$) admite al menos una combinación lineal estacionaria, condición necesaria para descartar una regresión espuria en niveles. El orden de rezagos del VAR subyacente se selecciona por mayoría de criterios de información (AIC, BIC, FPE, HQIC); en ninguna de las dos ventanas muestrales los cuatro criterios coinciden (en la ventana original seleccionan 2, 1, 2 y 2 rezagos respectivamente, y en la ampliada 5, 1, 5 y 2), de modo que la especificación final no se apoya en un único criterio sino que se reporta la sensibilidad del rango de cointegración y de la estimación posterior a esta elección (Sección \ref{sec:resultados_vecm}).

En la ventana ampliada, el contraste de Johansen exhibe además una sensibilidad sustantiva al punto de inicio de la muestra: con el rezago más parsimonioso (criterio BIC), la Traza no rechaza la ausencia de cointegración ($H_0: r=0$) al incluir el año 1999 completo, pero sí la rechaza al excluirlo. Esta fragilidad, documentada en el Capítulo \ref{ch:resultados}, es coherente con la literatura sobre la distorsión del contraste de Johansen ante quiebres estructurales severos no modelados en la tendencia determinística \citep{gregory1996,johansen2000}, y es precisamente el tipo de tensión que el pedido explícito del director de extender la ventana ``aun con cambios estructurales extremos'' buscaba hacer emerger: la muestra ampliada no ofrece una imagen más nítida de la cointegración de largo plazo, sino una menos concluyente en el margen, atribuible al colapso de la convertibilidad y la cesación de pagos de 2001--2002. Ante esta ambigüedad, se impone un rango de cointegración $r=1$ por motivo teórico (una única relación de largo plazo, la Función de Reacción Fiscal de \citet{bohn1998}), criterio ya utilizado en la ventana original ante una ambigüedad de signo opuesto (Sección \ref{sec:resultados_vecm}), documentando en ambos casos que la imposición no es una elección neutral.

\subsection{Etapa 3: Vectores con Corrección de Error (VECM) y VAR Restringido}
\label{sec:vecm}
A pedido del director de tesis, quien observó que la elección de Mínimos Cuadrados Ordinarios Dinámicos (DOLS) del trabajo original no había sido justificada frente a la alternativa de un sistema, la estimación de referencia de la Función de Reacción Fiscal de largo plazo pasa de DOLS a un Modelo de Vectores con Corrección de Error (VECM) sobre la ventana ampliada. La diferencia no es meramente técnica: DOLS estima la relación de largo plazo mediante una única ecuación, con $pb_t$ como variable dependiente y la endogeneidad de corto plazo tratada mediante adelantos y rezagos \textit{ad hoc} de $\Delta d_t$ \citep{stock1993}; un VECM, en cambio, modela el sistema completo, compuesto por $d_t$, $pb_t$, $risk_t$ y $tc_t$, como conjuntamente endógeno, sin necesidad de designar una variable dependiente ``principal'' ni de purgar sesgos con una construcción auxiliar: la endogeneidad queda integrada directamente en la estructura del sistema.

La especificación estimada es:
\begin{equation}
\Delta \mathbf{x}_t = \boldsymbol{\alpha} \boldsymbol{\beta}' \mathbf{x}_{t-1} + \sum_{i=1}^{k-1} \boldsymbol{\Gamma}_i \Delta \mathbf{x}_{t-i} + \boldsymbol{\delta}\, \tilde{y}_t + \boldsymbol{\varepsilon}_t
\end{equation}
donde $\mathbf{x}_t = (d_t, pb_t, risk_t, tc_t)'$, $\boldsymbol{\beta}$ es el vector de cointegración (relación de largo plazo, normalizado sobre $d_t$), $\boldsymbol{\alpha}$ contiene los coeficientes de ajuste (velocidad de corrección de cada variable hacia el equilibrio de largo plazo ante un desvío del período anterior) y la brecha del producto $\tilde{y}_t$ entra como variable exógena estacionaria, fuera del vector de cointegración, por ser I(0) por construcción. El rango de cointegración se impone en $r=1$ (Sección \ref{sec:etapa2_johansen}) y el rezago en primeras diferencias $k-1$ se fija según el criterio BIC sobre la ventana ampliada, reportando en el Capítulo \ref{ch:resultados} la sensibilidad de $\boldsymbol{\beta}$ y $\boldsymbol{\alpha}$ a esta elección.

\subsubsection{Identificación Estructural y VAR Restringido con Reglas Macroeconómicas}
\label{sec:svar_metodologia}
Para identificar la transmisión de perturbaciones macroeconómicas y fiscales sin recurrir a un ordenamiento recursivo de Cholesky arbitrario, se formula un Vector Autoregresivo Estructural Restringido (SVAR) siguiendo el esquema de identificación contemporánea de \citet{blanchard2002} y la restricción de acumulación de deuda de \citet{favero2007,favero2012}. Para el vector de cinco variables $\mathbf{Y}_t = (\tilde{y}_t, pb_t, \text{EMBI}_t, \text{TCRM}_t, d_t)'$, el sistema estructural $A_0 \mathbf{u}_t = B \mathbf{e}_t$ impone las siguientes restricciones basadas en teoría económica:
\begin{enumerate}
    \item \textbf{Rigidez de Decisión Presupuestaria}: El superávit primario reacciona contemporáneamente a la brecha del producto únicamente a través de la semi-elasticidad cíclica automática de la recaudación tributaria ($\alpha_y = 0{,}25$, estándar OCDE/FMI para Argentina; \citealt{girouard2005}; \citealt{daude2010}; \citealt{alberola2014}), exhibiendo rigidez institucional frente a shocks de riesgo soberano y tipo de cambio dentro del mismo trimestre ($a_{23} = a_{24} = 0$).
    \item \textbf{Ajuste Inmediato de Precios Financieros}: El riesgo soberano ($\text{EMBI}_t$) y el tipo de cambio real ($\text{TCRM}_t$) son variables de mercado financiero que absorben contemporáneamente las perturbaciones macroeconómicas y fiscales del período.
    \item \textbf{Identidad Dinámica de Acumulación}: La deuda pública evoluciona restringida por la restricción presupuestaria intertemporal:
    \begin{equation}
    d_t = \left( \frac{1 + r_t}{1 + g_t} \right) d_{t-1} - pb_t + sft_t
    \end{equation}
\end{enumerate}
A partir de la matriz de impacto estructural estimada $\hat{S} = A_0^{-1} B$, se calculan las Funciones de Impulso-Respuesta Estructurales (SVAR IRF) con bandas de confianza bootstrap al 95\% y la Descomposición de Varianza del Error de Pronóstico (FEVD).

DOLS no se descarta: se conserva como ejercicio de robustez y de comparación metodológica sobre la ventana original de 88 observaciones, con la especificación:
\label{sec:dols}
\begin{equation}
pb_t = \alpha + \rho\, d_{t-1} + \gamma\, \tilde{y}_t + \sum_{j=-4}^{4} \phi_j \Delta d_{t-j} + \varepsilon_t
\end{equation}
aumentada, siguiendo a \textcite{stock1993}, con $m=4$ adelantos y rezagos de la primera diferencia de la deuda para obtener un estimador superconsistente de $\rho$ aun ante incertidumbre sobre el rango exacto de cointegración del sistema.

\subsection{Etapa 4: Corrección por Endogeneidad (IV-2SLS)}
El tratamiento sistémico de la endogeneidad en el VECM (Etapa 3) no vuelve redundante este paso: IV-2SLS se aplica sobre la especificación DOLS de la ventana original, como corrección de endogeneidad dentro de un marco de ecuación única, complementaria y no sustituta del tratamiento estructural del VECM. Dado que la teoría económica predice causalidad inversa entre el resultado fiscal y el riesgo soberano, se instrumenta el EMBI+ ($risk_t$) mediante el índice de volatilidad global VIX y un diferencial de riesgo soberano regional (variable de aproximación de mercado construida sobre el ETF EMB, iShares JPMorgan USD Emerging Markets Bond, que replica el JPMorgan EMBI Global Diversified, la misma familia de índices que el EMBI+ argentino), estimando por Variables Instrumentales en Dos Etapas (IV-2SLS). La fortaleza del diseño instrumental se evalúa mediante el estadístico $F$ de la regresión de primera etapa (exigiendo $F > 10$, \citealp{staiger1997}), junto con el test de endogeneidad de Wu-Hausman y el test de sobreidentificación de Sargan.

\subsection{Etapa 5: Modelos de Umbral (Hansen y TVECM) y Fatiga Fiscal}
Para evaluar la existencia de fatiga fiscal, se estima una regla de reacción fiscal con efectos de umbral determinados endógenamente por el nivel de riesgo soberano, siguiendo la formulación de \citet{hansen1999}:
\begin{equation}
pb_t = \alpha + \beta_1 d_{t-1} \mathbb{I}(\text{EMBI}_t \le \tau) + \beta_2 d_{t-1} \mathbb{I}(\text{EMBI}_t > \tau) + \gamma \tilde{y}_t + \epsilon_t
\end{equation}
donde $\mathbb{I}(\cdot)$ es la función indicadora de régimen, $\tau^*$ denota el umbral crítico óptimo de riesgo soberano estimado mediante búsqueda de grilla sobre los percentiles 15--85 de la distribución empírica del EMBI+, y $\tilde{y}_t$ es la brecha del producto como variable de control macroeconómico. La significatividad estadística del quiebre se contrasta mediante el estadístico Sup-LM, cuya distribución bajo la hipótesis nula de linealidad se aproxima mediante re-muestreo con 1.000 réplicas \citep{hansen2000}.

Adicionalmente, se extiende el análisis al marco multivariado mediante el modelo \textbf{Threshold VECM (TVECM)} de \citet{hansen2002}, permitiendo que las velocidades de ajuste del vector $\boldsymbol{\alpha}$ conmuten de régimen según el estado del mercado financiero:
\begin{equation}
\Delta \mathbf{x}_t = \begin{cases}
\boldsymbol{\mu}_1 + \boldsymbol{\alpha}_1 \boldsymbol{\beta}' \mathbf{x}_{t-1} + \boldsymbol{\Gamma}_1 \Delta \mathbf{x}_{t-1} + \boldsymbol{\varepsilon}_t & \text{si } \text{EMBI}_{t-1} \le \tau^* \\
\boldsymbol{\mu}_2 + \boldsymbol{\alpha}_2 \boldsymbol{\beta}' \mathbf{x}_{t-1} + \boldsymbol{\Gamma}_2 \Delta \mathbf{x}_{t-1} + \boldsymbol{\varepsilon}_t & \text{si } \text{EMBI}_{t-1} > \tau^*
\end{cases}
\end{equation}

Sobre la ratio Deuda/PIB se implementó además el procedimiento formal de \textcite{bai2003}: programación dinámica exacta sobre la suma de cuadrados residuales de la serie, con selección endógena del número de quiebres mediante el criterio BIC \citep{yao1988}.

\subsection{Etapa 6: Simulación Estocástica y Análisis de Sostenibilidad de la Deuda (DSA)}
Para trascender las proyecciones deterministas convencionales, se implementa un Análisis de Sostenibilidad de la Deuda (DSA) estocástico siguiendo el estándar de organismos multilaterales \citep{imf2003,imf2021}. Se simulan escenarios deterministas y una simulación de Monte Carlo de 1.000 iteraciones, en la cual las variables macroeconómicas fundamentales reciben perturbaciones aleatorias extraídas de una distribución $t$ de Student multivariada con $\nu \approx 4.8$ grados de libertad acoplada a correlaciones dinámicas DCC-GARCH \citep{engle2002}.

\subsubsection{Modelización de la Tasa de Refinanciación mediante Proceso Cox-Ingersoll-Ross (CIR)}
\label{sec:cir_metodologia}
Para dotar a la tasa de interés de refinanciación externa ($r_{f,t} = r_{US,t} + \text{risk}_t$) de consistencia con la teoría de estructura temporal, se modela el riesgo soberano como un proceso de difusión continua con reversión a la media \citep{cox1985}:
\begin{equation}
d(\text{risk}_t) = \kappa (\theta - \text{risk}_t) dt + \sigma \sqrt{\text{risk}_t} \, dW_t
\end{equation}
donde $\kappa$ es la velocidad de reversión a la media, $\theta$ es el nivel de equilibrio de largo plazo y $\sigma$ es el parámetro de volatilidad. La especificación garantiza la no negatividad estricta del riesgo país siempre que se satisfaga la condición de Feller:
\begin{equation}
2\kappa\theta > \sigma^2
\end{equation}
La estimación se realiza por Máxima Verosimilitud exacta mediante la función de densidad de transición Chi-cuadrado no central expresada en términos de funciones modificadas de Bessel de primera especie \citep{cox1985,duffie1996}, integrando las trayectorias simuladas directamente en la ley de movimiento del ratio de deuda.
